\begin{titlepage}
\centering
\vspace*{14mm}
{\sffamily\small 系统学习教材\par}
\vspace{13mm}
{\sffamily\bfseries\fontsize{29}{40}\selectfont 机械臂控制与仿真\par}
\vspace{3mm}
{\sffamily\bfseries\fontsize{25}{35}\selectfont 详细教材\par}
\vspace{11mm}
{\sffamily\Large 从数学推导到七轴柔顺控制\par}
\vspace{7mm}
{\large 概念解释\quad 完整数值算例\quad 公式推导\quad 精确图解\par}
\vspace{17mm}
\includegraphics[width=.86\linewidth]{figures/rt_cartesian.pdf}
\vfill
\begin{minipage}{.84\linewidth}
\normalsize
以单关节和两连杆为计算主线，连接七轴机械臂的参数辨识、低阻抗控制、零空间、示教回放与相机云台。

本书根据两份上传教材重新整合编写，统一符号与假设，补充中间计算、算法边界和可复现图形。
\end{minipage}
\vspace{13mm}

{\small 整合与扩展版\quad 2026年11月23日\par}
\end{titlepage}

\chapter*{本书的目的与阅读方法}
\addcontentsline{toc}{chapter}{本书的目的与阅读方法}

这本书面向已经接触机械臂软件、希望把公式和工程现象连起来的读者。学习目标是：看到机械臂的位置、速度、力矩、外力和误差时，知道这些量分别来自哪里，能把一条控制公式拆成可解释、可计算、可检验的步骤。

两份原稿包含很广的知识，但数学、控制和工程内容反复交叉，部分独立公式仍显示为源码，部分图只有标题，某些结论缺少成立条件。新版本重新组织知识顺序，并对符号、数值和前提进行统一；它不是把两个 PDF 顺次拼接，也不是删减成知识点摘要。

\section*{原始材料和引用约定}

\begin{description}[leftmargin=2.6em,style=nextline]
\item[原稿 A]《机械臂与机器人控制算法 从零开始的基础巩固手册》，上传文件共90页。侧重数学基础、经典与现代控制、机器人控制和工程能力体系。
\item[原稿 B]《机械臂与仿真工程教材 数学及格也能懂版》，上传文件共129页。侧重运动学、辨识、摩擦、仿真、示教、遥操作和七轴云台。
\end{description}

每章末标注原稿对应页码，页码均指上传 PDF 的物理页序。涉及补充的一手软件文档和论文，采用参考文献编号。网页资料核对日期为2026年11月23日；软件示例应对应实际安装版本。原稿提到但本次没有提供完整内容的其它讲义或教程，不会被当成已经通读的独立来源。

\section*{怎样读一个公式}

每遇到一个公式，依次问五件事：它在解决什么问题；每个符号是什么量、在哪个坐标系、什么单位；等号两边的维度是否相同；它由哪一步物理关系或数学运算得到；它在什么前提下可用。只记最后一行，通常不足以排查真实代码。

例如 $\tau=J^\mathsf TF$ 并不只是“雅可比转置乘力”。必须知道 $J$ 映射哪个参考点的速度，$F$ 是哪一方作用于哪一方，力和速度是否在同一个坐标系，以及六维向量采用哪种行顺序。正确的单位、坐标与功率关系，可以比盲调参数更早发现错误。

\section*{本书每章的组织方式}

正文先说明问题和直觉，再引入概念与符号，随后给出推导与逐步计算，图形用来显示公式的结果或边界，最后用计算题和解答检查掌握程度。解答紧跟本章，便于自学；建议先遮住解答手算，再逐步对照。

全书包括23章正文、4个附录、57幅图、114个框内算例，以及83道附逐步解答的章末计算练习；正文中还包含没有单独入框的数值演算。目录和图目录可点击跳转，文末列出原稿与补充的一手参考资料。

所有新生成的曲线均来自明确方程或给定数据。示意几何和合成仿真会在图注中说明，不会写成真机结果。章节中出现的教学刚度、阻尼、采样周期和质量，只用于解释计算过程；迁移到 NERO 需要代入实际模型、接口和测量。

\section*{知识地图}

\begin{longtable}{p{24mm}p{44mm}p{86mm}}
\toprule 章节 & 主要问题 & 学完后应能亲手完成的计算\\\midrule
\endhead
第1--2章 & 数学工具 & 解线性方程，解释秩与条件数，求导积分，解简单动态方程，理解变换与离散化\\
第3--5章 & 位姿与运动学 & 坐标转换，旋转误差，2R正逆解，雅可比、奇异与速度零空间\\
第6--7章 & 动力学与执行器 & 从能量推导力矩，理解惯量耦合、电机、传动、柔性和接触\\
第8--10章 & 参数辨识与标定 & 建回归矩阵，算加权最小二乘，设计激励，分析摩擦、负载与手眼标定\\
第11--14章 & 控制与估计 & 设计PD和PID，分析频域，算LQR和离散模型，逐步更新滤波器与观测器\\
第15--16章 & 力控与七轴柔顺 & 推导阻抗、导纳和操作空间方程，检查零空间力矩何时不干扰主任务\\
第17--19章 & 运动与人机输入 & 解五次轨迹与S曲线，处理示教，计算DMP、ProMP、遥操作和视觉映射\\
第20--21章 & 工程和仿真 & 对齐接口与时序，估总线带宽，建立多速率仿真，区分模型真值与测量\\
第22--23章 & 约束和综合应用 & 手算MPC与任务权重，理解WBC，把整条链迁移到七轴云台\\
\bottomrule
\end{longtable}

\section*{按当前问题选择入口}

若当前最困惑的是“参数辨识到底在算什么”，先读第1章线性代数与最小二乘，再读第6、8、9章。若重点是“笛卡尔拖动和肘部都要软”，先读第3、5、6章，再读第15、16章。若要做运镜与示教，连接第17--21章及第23章。跳读时遇到尚不理解的符号，可先查附录的术语表，再回到首次推导的章节。

不按固定周数承诺学习速度。一个可靠的前进标准是：能不看答案重算本章的关键数字，并说清楚改一个参数后，图形会怎样变化。对于本书的二维或一维例子，这比单纯背诵算法名称更有帮助。

\chapter*{全书符号和三个贯穿模型}
\addcontentsline{toc}{chapter}{全书符号和三个贯穿模型}

\section*{全书默认约定}

关节角用弧度，时间用秒，长度用米，力用牛顿，转动力矩用牛米。角度虽然在 SI 量纲中可视为无量纲，工程表达仍保留 rad，以便区分角位移、线位移及对应增益。

\begin{longtable}{p{33mm}p{79mm}p{42mm}}
\toprule 符号 & 含义 & 常见单位或维度\\\midrule\endhead
$q,\dot q,\ddot q$ & 关节位置、速度、加速度列向量 & rad，rad/s，rad/s$^2$\\
$q_d$ 与 $e=q-q_d$ & 期望关节角与实际减期望误差 & rad\\
$\tau_{\rm cmd},\tau_{\rm ext}$ & 执行器关节侧命令，环境对机器人的外力矩 & N m\\
$M(q)$ & 关节空间惯量矩阵 & kg m$^2$，$n\times n$\\
$C(q,\dot q)\dot q$ & 科氏与离心合成项 & N m\\
$g(q),\tau_f$ & 重力项与左端约定的摩擦项 & N m\\
$p,R,T$ & 位置、旋转矩阵、齐次变换 & $3\times1$，$3\times3$，$4\times4$\\
$R_{AB}$ & 把B系向量分量转换到A系 & $v_A=R_{AB}v_B$\\
$V=(v^\mathsf T,\omega^\mathsf T)^\mathsf T$ & 线速度在前，角速度在后 & m/s，rad/s\\
$F=(f^\mathsf T,m^\mathsf T)^\mathsf T$ & 力在前，力矩在后；与$V$同点同轴 & N，N m\\
$J$ & 对应任务的几何或明确指定的误差雅可比 & 行数由任务决定\\
$H,\beta,y$ & 堆叠回归矩阵、参数、测量目标 & 由回归定义\\
$T_s,h$ & 控制采样周期、物理积分步长 & s\\
$K,D$ & 刚度与阻尼；按关节或任务空间说明 & 由空间和物理量决定\\
\bottomrule
\end{longtable}

统一的动力学符号为
\begin{equation*}
M(q)\ddot q+C(q,\dot q)\dot q+g(q)+\tau_f
=\tau_{\rm cmd}+\tau_{\rm ext}.
\end{equation*}
因此，采用 $e=q-q_d$ 的恢复反馈力矩一般带负号。个别小节为了连接标准优化或经典控制推导采用参考减输出，会在当地明确定义，不能跨小节直接移植符号。

\section*{模型一 单关节}

惯量 $J_r=0.1\,\mathrm{kg\,m^2}$，黏性系数 $b=0.2\,\mathrm{N\,m\,s/rad}$，基础方程 $J_r\ddot q+b\dot q=\tau$。它用于把控制、滤波和离散化算清楚。$J_r$ 是单轴转动惯量，与几何雅可比 $J(q)$ 不同。

\section*{模型二 竖直平面两连杆}

$l_1=l_2=0.5\,\mathrm m$，$m_1=m_2=1\,\mathrm{kg}$，$c_1=c_2=0.25\,\mathrm m$，绕运动平面法向的质心惯量 $I_1=I_2=1/48\,\mathrm{kg\,m^2}$。重力沿 $-y$，$q_1$ 从水平正 $x$ 轴起算，$q_2$ 是第二杆相对第一杆的角。少数独立算例采用其它数值时会明确重列参数。

\section*{模型三 七轴机械臂与相机任务}

以固定基座七轴机械臂作应用背景，但不虚构真实 DH、URDF 惯量、驱动器力矩限制或厂商内部算法。其主要用途是解释六维任务、冗余、零空间、工具负载和世界到相机的误差链。物理图和低维算例用于说明原理，不替代 NERO 实际模型。

\clearpage
\tableofcontents
\clearpage
\listoffigures
\clearpage
